\section{Exact terminal attachment}\label{sec:terminal}
For a critical product $D$, let $N_D(T)$ count subgroups of $D\times T$
full on $T$ and on each displayed critical factor. The group $T$ is
the complete exterior projection, with all of its internal couplings.
Write $d_2(T)=\dim\Hom(T,C_2)$ and
$A_2(T)=T/(T'T^2)$.

\begin{lemma}[Binary contraction]\label{lem:contraction}
If $O^2(T)$ is the intersection of the normal subgroups of $T$ with
$2$-group quotient, image and full preimage induce a bijection
\[
 N_D(T)=N_D(T/O^2(T)).
\]
\end{lemma}
\begin{proof}
Goursat's lemma expresses a subgroup above its two projections by a
common quotient. The projection inside $D$ is a $2$-group, so the common
quotient of $T$ is a $2$-group. Its kernel contains $O^2(T)$.
Every subgroup being counted therefore contains $1\times O^2(T)$.
\end{proof}

\begin{lemma}[Central extension fibre]\label{lem:central-fibre}
Let $1\to K\to X\to Y\to1$ be central with $K$ elementary abelian
of exponent two, and let $\xi\in H^2(Y,K)$ be its class. Define
\[
 \Ann(\xi)=\{\lambda\in K^*: \lambda_*\xi=0\text{ in }H^2(Y,\F_2)\}.
\]
The number of subgroups of $X$ mapping onto $Y$ equals
\begin{equation}\label{eq:central-fibre}
 \sum_{L\leq\Ann(\xi)}2^{d_2(Y)\dim L}.
\end{equation}
\end{lemma}
\begin{proof}
The annihilator is the image of restriction $\Hom(X,C_2)\to K^*$:
a normalized section identifies extension of a character with its scalar
defect being a coboundary. Fix $W=H\cap K$. Taking kernels bijects these
subgroups with homomorphisms $r:X\to K/W$ restricting to the quotient
map on $K$. Such an $r$ exists exactly when every scalar coordinate extends,
or $W^\perp\leq\Ann(\xi)$. Differences of extensions factor through $Y$,
so this fibre is a torsor under $\Hom(Y,K/W)$, of size
$2^{d_2(Y)\codim W}$. Different intersections are disjoint.
\end{proof}

Let $V=D/K=\F_2^r$, $c=\dim K$ and $\delta=r/2-c$. For an arbitrary
finite group $T$ define
\begin{equation}\label{eq:tau}
 d=d_2(T),\qquad
 \tau=\dim\ker\bigl(H^2(A_2(T),\F_2)\longrightarrow H^2(T,\F_2)\bigr),
 \qquad X_r(d)=\sum_j\gauss rj2\,2^{dj}.
\end{equation}
The kernel in this definition is retained as a subspace, not merely
replaced by the number $\tau$, until a positive upper bound is taken.

\begin{proposition}[Exact quotient records]\label{prop:terminal-record}
For $0\leq u\leq r$, let $\mathcal P_u$ be the linear maps
$p:\F_2^u\oplus A_2(T)\to V$ injective on the first summand and
surjective onto each displayed critical quotient coordinate. Pull back
the critical extension to $C_2^u\times T$, and call its class $\xi_p$.
Then
\begin{equation}\label{eq:terminal-record}
 N_D(T)=\sum_{u=0}^r\frac{1}{|\GL(u,2)|2^{ud}}
 \sum_{p\in\mathcal P_u}\ \sum_{L\leq\Ann(\xi_p)}2^{(u+d)\dim L}.
\end{equation}
\end{proposition}
\begin{proof}
The image $Y\leq V\times T$ of a subgroup being counted is onto $T$.
Put $U=Y\cap V$. Its graph map $T\to V/U$ factors through $A_2(T)$.
A linear lift of this map identifies $Y$ with $U\times T$.
There are exactly $|\GL(u,2)|2^{ud}$ representations by a choice of
ordered basis for $U$ and a linear lift. Lemma~\ref{lem:central-fibre}
counts every subgroup above $Y$; surjectivity modulo the critical
Frattini kernels gives the required full projections.
\end{proof}

\begin{proposition}[Uniform terminal bound]\label{prop:terminal-bound}
With the notation above,
\begin{equation}\label{eq:terminal-double}
 N_D(T)\leq\phi^{-2}\sum_{j=0}^r\sum_{\ell=0}^c
 \gauss c\ell2\,72^\ell 2^{\tau\ell}
                    2^{(j-\ell)(r-j+d)}.
\end{equation}
Consequently
\begin{equation}\label{eq:terminal-ratio}
 \frac{N_D(T)}{X_r(d)}\leq (r+2)^{O(1)}
 \left(1+\sum_{\ell\geq1}
 2^{-\ell(\delta+d/2-\tau-\log72)-3\ell^2/4}\right).
\end{equation}
In particular, if $\tau=0$ then
$N_D(T)\leq C(r+2)^kX_r(d)$ for absolute $C,k$.
\end{proposition}
\begin{proof}
Put $u=r-j$ and $k=u+d$. For an $\ell$-dimensional space of splitting
characters in $K^*$ use row-echelon coordinates as in
Proposition~\ref{prop:extra-lifts}. The original product section has a
bilinear scalar defect with diagonal $q_\lambda$. If its pullback to
$U\times T$ is a coboundary, evaluation on a vertical involution,
comparison on commuting vertical and exterior elements, and restriction
to $T$ give, respectively,
\[
 q_\lambda|_U=0,\qquad
 \beta_\lambda(U,f(A_2(T)))=0,\qquad
 q_\lambda\circ f\in\ker\bigl(H^2(A_2(T))\to H^2(T)\bigr).
\]
Diagonal evaluation $c(a,a)-c(0,0)$ descends to a linear map on actual
$H^2(A_2(T),\F_2)$, since it vanishes on coboundaries. The first two
conditions force the whole quadratic outcome to be the exterior pullback
of the diagonal of a class in the retained inflation kernel. This allowed
function space has dimension at most $\tau$. Modulo that space, the
joint row relations determine every pivot outcome. Each outcome has at
most $2^\tau$ representatives, and Lemma~\ref{lem:quadratic} leaves at
most $72$ surjective maps for each representative on the whole
$k$-dimensional domain.
Dropping the other restrictions leaves at most
$72^\ell2^{\tau\ell}2^{(r-2\ell)k}$ ordered maps.
The record divisor is at least $\phi2^{uk}$. Multiplication by the
fibre size $2^{k\ell}$ and summation over $L$ prove
\eqref{eq:terminal-double}, even with one fewer factor $\phi^{-1}$.

Use $\gauss c\ell2\leq\phi^{-1}2^{\ell(c-\ell)}$ and write
$A=r+d$. The exponent apart from this fixed constant is
\[
 j(A-j)+\ell(c+\tau+\log72-\ell-A+j).
\]
If $d\leq r$, its unconstrained maximum in $j$ is
\[
 A^2/4-\ell(\delta+d/2-\tau-\log72)-3\ell^2/4,
\]
whereas an integer $j$ nearest $A/2$ gives
$X_r(d)\geq2^{A^2/4-O(1)}$. If $d\geq r$, the exponent is increasing
on $0\leq j\leq r$; its maximum is at $j=r$. Divide by the literal
term $2^{rd}\leq X_r(d)$. The remaining exponent is
$\ell(r/2-\delta+\tau+\log72-d-\ell)$, which is smaller than the
one in \eqref{eq:terminal-ratio} by
$\ell(d-r)/2+\ell^2/4$. The $r+1$ choices of $j$ are a polynomial
factor. Finally $\delta,d\geq0$, so at $\tau=0$ the series is bounded
by $\sum_{\ell\geq0}72^\ell2^{-3\ell^2/4}<\infty$.
\end{proof}

\begin{remark}
The module-extension annihilator used below is different from the
group-extension annihilator in \eqref{eq:central-fibre}. No argument
identifies the two. In particular a scalar quadratic form which is
nonzero on $A_2(T)$ may become zero after inflation to $T$.
\end{remark}
